% slides setup
\documentclass{beamer}
\usecolortheme{norcdark}

% handouts setup
% \documentclass[handout]{beamer}
% \usepackage{pgfpages}
% \pgfpagesuselayout{2 on 1}[letter,border shrink=5mm]
% \usecolortheme{norclight}

\usetheme{boadilla}

\beamertemplatenavigationsymbolsempty

\usepackage{amsmath,amsfonts,natbib,comment,graphicx,roboto}

\newcommand{\Expect}{\mathbb{E}}
\newcommand{\Var}{\mathbb{V}}
\newcommand{\One}{{\rm 1\kern-0.3em I}}

\begin{document}

\title{\textbf{Analysis of Complex Health Survey Data}}
\author[Stas Kolenikov]{Stas Kolenikov \\ @StatStas}
% \date{ICHPS 2023 \\ Conference WiFi: ?? \\ GitHub: https://github.com/skolenik/ICHPS2023-svy }
\date[ICHPS 2023]{ICHPS 2023
 % \medskip \\ Conference WiFi: Marriott{\_}CONFERENCE; CHSBQT
 \medskip \\ GitHub: https://github.com/skolenik/ICHPS2023-svy }

\institute[NORC]{NORC \\ \medskip Views and opinions are those of the instructor only, 
    \\ not those of NORC or ASA}

% redefine beamer colors

\logo{\includegraphics[height=0.4cm]{NORCLogo_Short_Color}}

\begin{frame}
    \titlepage
\end{frame}

\begin{frame}\frametitle{Target audience and objectives}

This short course is aimed at the audience with general knowledge of statistics
at graduate level. No knowledge of survey statistics is assumed
\begin{itemize}
    \item but a 2 hour course won't make you an expert in it, either!
\end{itemize}

The course will cover:
\begin{itemize}
    \item Role of survey data in today's multi-faceted world of data collection
    \item Basic features of survey data
    \item Basic principles of survey inference
    \item The most important peculiarities of analyzing survey data
\end{itemize}

By the end of the course, you will be better prepared to read technical reports
and papers describing complex health survey data. You will be better equipped to
understand the syntax of the complex health survey data analysis modules
in R and Stata (self-study).

\end{frame}

\begin{frame}\frametitle{What informs health policy?}

If we want to improve health outcomes of constituents of health policy,
how are the differences in health outcomes related to
the differences in what may potentially affect health?

\begin{itemize}
    \item Biological pathways and experiments (pharma RCT)
    \item Biological pathways and observational studies
    \begin{itemize}
        \item leading to universally applicable health policies
        \item e.g., smoking and tobacco
    \end{itemize}
    \item Socioeconomic determinants and observational studies (survey data)
    \begin{itemize}
        \item leading to allocation of limited resources across population groups
        \item e.g., targeted cancer screening or immunization campaigns
    \end{itemize}
\end{itemize}

\end{frame}

\section{Examples of complex survey data}
% 1. Examples of complex health survey data: NHIS, NHANES, NIS, BRFSS. Survey design trade-offs: frames, coverage, modes and costs. Public use and restricted data sets.

\begin{frame}\frametitle{What is a complex survey?}

If the methodology description of your data set contains keywords or phrases like
\begin{itemize}
    \item stratified sampling
    \item multistage sampling
    \item random digit dialing
    \item nonresponse adjustment
    \item probability weights
    \item replicate weights
\end{itemize}

\ldots it is a complex survey data set, and you are in the right workshop to learn more about them!

\medskip

\citet{korn:graubard:1999}, \citet{heeringa:west:berglund:2017}, \citet{lumley:2010},
\citet{chambers:skinner:2003}, \citet{pfeffermann:rao:2009}, \citet{wu:thompson:2020}.

\end{frame}

\begin{frame}\frametitle{Finite population surveys}

\begin{itemize}
    \item Developed initially (1920--1930s) to study finite populations with fixed characteristics
        \begin{itemize}
            \item volume of a tree
            \item last year revenue of a business
            \item total grain harvest by a farm
            \item person's sex
        \end{itemize}
    \item Random variables are 0/1 indicators of being in the sample vs. not
    \item Inference/statistical distributions are with respect to the sampling mechanism
    \item 1970s and on: statistical models are incorporated and blended with sampling inference
\end{itemize}

\end{frame}

\begin{frame}\frametitle{Sample surveys vs. anything else}

\begin{itemize}
    \item Clinical trials
        \begin{itemize}
            \item assume biology is the same, hence sample selection is not an issue
            \item inference is wrt randomized assignment
        \end{itemize}
    \item Administrative records
        \begin{itemize}
            \item by-products of systems built by somebody else
            \item created for reasons other than research (e.g., billing)
            \item concepts/definitions/categories may not match the research questions
            \item cannot uncover how people adapt to the admin structures
        \end{itemize}
    \item Big Data
        \begin{itemize}
            \item terrific at answering the ``what is happening'' question
            \item sweeps rationale/motives for human behavior under the carpet --- awful at ``why'' question
            \item ML/AI methods have $\approx$ zero transparency
            \item ``$N=$ all'' is a misleading slogan
        \end{itemize}
\end{itemize}

\end{frame}

\begin{frame}\frametitle{Dimensions to break down}

\begin{itemize}
    \item Human populations vs. establishment populations (hospitals, practices, providers)
    \item Sampling frames and modes
    \begin{itemize}
        \item List of patients
        \item Population register
        \item Area sample of the general population : face-to-face/specimen interview
        \item Random digit dialing sample of the general population : phone interview
        \item Address-based sample of the general population : mail or face-to-face interview
        \item Respondent driven samples : SAQ or face-to-face interview
        \item Screening for target population
        \item Combinations: multiple frame and/or multiple mode surveys
    \end{itemize}
    \item Probability panels
\end{itemize}

\end{frame}

\begin{frame}\frametitle{Health surveys in the U.S.}

\begin{itemize}
    \item National Health Interview Survey (NHIS): area sample, face-to-face; CDC/NCHS
    \item National Health And Nutrition Examination Survey (NHANES): area sample, face-to-face, biological specimen; CDC/NCHS
    \item Behavioral Risk Factor Surveillance Survey (BRFSS): phone survey; CDC, state health agencies
    \item National Survey of Drug Use and Health (NSDUH): area sample, face-to-face; SAMHSA
    \item Consumer Assessment of Healthcare Providers and Systems (CAHPS): list sample of patients, multimode; AHRQ+CMS
    \item Medical Expenditure Panel Survey (MEPS): three-fold survey of households, their insurers, and their medical providers; AHRQ
    \item National Inpatient Sample of the Healthcare Cost and Utilization Project (HCUP-NIS): 
    	sample of discharges; discharge and hospital data elements; AHRQ
\end{itemize}

\end{frame}

\begin{frame}\frametitle{International health surveys}

\begin{itemize}
    \item Demographic and Health Surveys (DHS) program:
    \begin{itemize}
        \item reproductive age women
        \item area samples (cluster $\approx$ village), face-to-face
        \item $\sim$90 countries, every $\sim$3--8 years
    \end{itemize}
    \item Multiple Indicator Cluster Surveys (UNICEF):
    \begin{itemize}
        \item women and children
        \item area samples (cluster $\approx$ village), face-to-face
        \item $\sim$120 countries, irregular $\sim$3--15 years
    \end{itemize}
    \item International Tobacco Control Policy Evaluation Project: sample design varies by country
    \begin{itemize}
        \item phone and web in developed countries
        \item anything goes in developing countries, face-to-face augmented with other modes
    \end{itemize}
\end{itemize}

\end{frame}

\section{Sampling designs and complex survey features}
% 2. Features of complex surveys: weights, clusters, strata, nonresponse adjustments.

\subsection{Terminology}

\begin{frame}\frametitle{Language}

\textit{Population} $\mathcal U$ $\equiv$ units for which the generalizable knowledge is sought; $|\mathcal{U}|=N$.

\smallskip

\textit{Observation unit} $\equiv$ the entity on which the survey measurements are
taken.

\smallskip

\textit{Sample} $\mathcal S$ $\equiv$ units that are selected for observation; $|\mathcal{S}|=n$.

\smallskip

\textit{Frame} $\equiv$ a method to identify, and often contact, a unit from the target population; a link between the target and feasible population.

\smallskip

\textit{Coverage} $\equiv$ relation between frame(s) and population.
A lot of the times, no single frame covers the entire population of interest.

\smallskip

\textit{Sampling unit} $\equiv$ the entity obtained from the sampling frame
with a single draw;
may match the observation units, or be a group of observation units.

\smallskip

\textit{EPSEM design} $\equiv$ equal probability of selection method; a sampling
design in which all observation units have the same probability of selection.
\textbf{Does not equate i.i.d. sample.}

\end{frame}

\subsection{Big Four: stratification, clustering, prob selection, weighting}

\begin{frame}\frametitle{Features of complex surveys}

The Big Four features are:

\begin{itemize}
    \item Stratification
    \item Clustering
    \item Unequal probabilities of selection
    \item Weight adjustments
\end{itemize}

Other statistical features may include
\begin{itemize}
    \item multiple phase sampling
        \begin{itemize} \item subsampling respondents for specimen data collection \end{itemize}
    \item multiple frames
        \begin{itemize} \item all phone surveys use both landline and cell frames \end{itemize}
    \item mode effect adjustments (prominent in CAHPS)
\end{itemize}

\end{frame}

\begin{frame}\frametitle{Stratification}

Stratification $\equiv$ breaking up the population/frames into mutually exclusive groups before sampling.
\begin{itemize}
    \item Geographic regions in f2f samples
    \item Diagnostic groups in patient list samples
\end{itemize}

Why?

\begin{itemize}
    \item Oversample subpopulations of interest if they can be identified on the frame(s)
    \item Oversample areas of higher concentration of the target rare population
    \item Ensure specific accuracy targets in subpopulations of interest
    \item Utilize different sampling designs/frames in different strata
    \item Balance things around/avoid weird outlying samples/spread the sample across the whole population
    \begin{itemize} \item Deeply stratified samples: dozens/hundreds of strata, 2 PSUs per stratum \end{itemize}
\end{itemize}



\end{frame}

\begin{frame}\frametitle{Cluster samples}

Cluster, or multistage sampling design $\equiv$ sampling groups of units rather than the ultimate observation units.
\begin{itemize}
    \item Geographic units (e.g., census tracts) in f2f samples
    \item Entities in natural hierarchies (e.g., hospitals/practices and providers within a practice)
\end{itemize}

Why?

\begin{itemize}
    \item Lists of observation units are not available, but survey statistician
        can obtain lists of higher level units
        for which residence or health service of the ultimate observation units
        can be established
    \item Reduce interviewer travel time/cost in f2f surveys
    \item Interest in multilevel modeling of hierarchical structures
\end{itemize}

Terminology: PSU $\equiv$ primary sampling unit $\approx$ cluster

\end{frame}

\begin{frame}\frametitle{Unequal probabilities of selection}

Why?

\begin{itemize}
    \item Oversample (smaller) subpopulations of interest (e.g., ethnic/racial minorities)
    \item Oversample areas of higher concentration of the target rare population
    \item Result of multiple stage/cluster sampling
    \begin{itemize}
        \item Most f2f samples are drawn using probability proportional to size (PPS) 
            sampling at the first few stages,
            fixed sample size at last stage $\Rightarrow$ approximately EPSEM
        \item If measures of sizes are not accurate, or differential nonresponse is 
            encountered, no longer EPSEM
    \end{itemize}
    \item Unintended result of multiple frame sampling
        \begin{itemize}
            \item dual phone users, i.e., those who have both landline and cell phone service, are more likely to be selected
        \end{itemize}
\end{itemize}

\end{frame}

\begin{frame}\frametitle{Weight adjustments}

Why? Corrections for\ldots

\begin{itemize}
    \item eligibility
    \item nonresponse (unavoidable in real world)
    \item frame noncoverage
    \item frame overlap in multiple frame surveys
    \item statistical efficiency
\end{itemize}

\citet{kalton:flores:cervantes:2003},
\citet{valliant:dever:kreuter:2013},
\citet{valliant:dever:2017},
\citet{kolenikov:2016}

\end{frame}

\begin{frame}\frametitle{Sampling is about doing the best job for the \$\$\$}

In the end of the day, all of the sampling features are there for 1+ of the following reasons:

\begin{itemize}
    \item Save money
        \begin{itemize}
            \item use cluster samples to save on travel costs
            \item use stratified samples to realize statistical efficiency gains
        \end{itemize}
    \item Cannot get the full population listing
        \begin{itemize}
            \item \ldots so have to use area samples to gradually zoom down to individuals
            \item \ldots so have to use infrastructure created for a different purpose (telecom or postal) to contact people
        \end{itemize}
    \item Overcome the real world data collection difficulties
        \begin{itemize}
            \item nonresponse weight adjustments
        \end{itemize}
\end{itemize}

\end{frame}

\section{Impact of complex survey features on analysis}
% 3. Impacts of weights, clusters, stratification on estimates and standard errors: descriptive statistics and regression models. Design effects.

\subsection{Biases}

\begin{frame}\frametitle{Design-based inference}

Recall the finite population sampling paradigm:
\begin{itemize}
    \item observed characteristics are fixed (e.g., person's height and weight, presence of a medical condition)
    \item random variables are 0/1 indicators of being vs. not being in the sample
\end{itemize}

Remember $\langle \Omega, \mathcal{F}, P \rangle$ from your Billingsley-based class?
\begin{itemize}
    \item elementary outcome = particular sample, i.e., subset of the population $\mathcal U$
    \item probability associated with the outcome = probability to draw a particular sample
    \item Hence $\Omega$ is a discrete space (although combinatorially large)
    \item \ldots all events are finite unions of elementary outcomes, and
    \item \ldots all sampling distributions are histograms
\end{itemize}

\end{frame}

\begin{frame}\frametitle{Design-based inference, cont.}

``Easy'' statistics: those linear in random variables.

\bigskip

We like linear statistics because expectations can be carried through,
and we can produce CLT-type results through multiplication of characteristic functions.

\begin{itemize}
    \item Trick question: is the sample mean $\bar y = \frac 1n\sum_{i=1}^n y_i$ a linear statistic?
	 \item Reminder: survey statisticians think in terms of sample inclusion indicators 
	       $\One_i, i\in\mathcal{U}$.
    \item Recast the sample mean in those terms: 
          $\bar y = \frac{\sum_{i\in\mathcal{U}} y_i \One_i}{\sum_{i\in\mathcal{U}} \One_i}$,
          which is a ratio of two random variables
	 \begin{itemize}
	      \item \ldots unless the sample size $n$ is fixed by design, which rarely happens in practice
	 \end{itemize}
\end{itemize}

\end{frame}

\begin{frame}\frametitle{Design-based inference, cont.$^2$}

So, what \textit{are} the linear functions of random variables? In survey statistics, these are the totals.

\begin{itemize}
    \item Target of inference (unobserved): population total $T[y] = \sum_{i\in\mathcal{U}} y_i$
    \item Sample totals (with weights): $t_w[y] = \sum_{i\in\mathcal{S}} w_i y_i \equiv \sum_{i\in\mathcal{U}} \One_i w_i y_i$
\end{itemize}

How to unbiasedly estimate the population total?
\begin{itemize}
    \item Define probability of selection ${\rm Prob}[i\in\mathcal{S}]=\pi_i$, then
    \item $\Expect_p t_w[y] = \sum_{i\in\mathcal{U}} \One_i w_i y_i = \sum_{i\in\mathcal{U}} \pi_i w_i y_i$
    \item To make this unbiased for $T[y]$, irrespective of values of $\{ y_i| i\in \mathcal{U}\}$,
        it makes sense to set weight so that $\pi_i w_i = 1 \, \forall i$, i.e. $w_i = 1/\pi_i$.
\end{itemize}

We just derived the \citet{horvitz:thompson:1952} estimator!

\bigskip

Note that we made no assumptions regarding $y_i \, \Rightarrow$

\centering

\textsc{survey inference is essentially nonparametric}.

\end{frame}

\begin{frame}\frametitle{Design-based inference, cont.$^3$}

What can we say about the means of sampling distributions (point estimation)?

\begin{itemize}
    \item To produce unbiased estimates of totals, we definitely need the Horvitz-Thompson weights.
    \item All other statistics are nonlinear, e.g., the weighted mean is a ratio
        $$ \bar y_w = \frac{\sum_{i\in\mathcal{S}} w_i y_i}{\sum_{i\in\mathcal{S}} w_i} \equiv \frac{t_w[y]}{t_w[1]} $$
    \item If some sort of LLN holds, then functions of consistently estimated quantities
        will be consistent, so\ldots
    \item \ldots it appears to make sense to use weights for all analyses.
\end{itemize}

\end{frame}

\begin{frame}\frametitle{Design-based inference, cont.$^4$}

Weighted mean again:
$$ \bar y_w = \frac{\sum_{i\in\mathcal{S}} w_i y_i}{\sum_{i\in\mathcal{S}} w_i} \equiv \frac{t_w[y]}{t_w[1]} $$

\bigskip

From the expression for $\bar y_w$, it can be seen that weights ``may not matter'' 
if $y_i$ are ``uncorrelated'' with weights $w_i$\ldots
\begin{itemize}
    \item \ldots although the sense of ``uncorrelated'' is very poorly defined
\end{itemize}
Typically, survey statisticians exploit knowledge of the features of the population of interest
in their sampling designs, so expecting that measurements and weights would be ``uncorrelated'' is
unwarranted, if not dangerous, for producing unbiased estimates.

\end{frame}

\begin{frame}\frametitle{Design-based inference, cont.$^5$}

What can we say about the variances of sampling distributions (standard errors)?

\begin{itemize}
    \item Variances are associated with second moments of random variables
    \item The random variables in surveys are $\One_i$, so\ldots
        \begin{itemize}
            \item $\Expect \One_i^2 = \Expect \One_i = \pi_i$
            \item $\Expect \One_i \One_j \equiv \pi_{ij} = {\rm Prob}[\mbox{both }i,j \in \mathcal{S}]$
        \end{itemize}
    \item Thus, sampling variances depend on the \textit{first order}
        probabilities of selection, and especially on the
        \textit{joint}, or \textit{second order} probabilities of selection:
\end{itemize}

$$
\Var( t[y] ) = \sum_{i, j\in\mathcal{U}} \frac{\pi_{ij}-\pi_i \pi_j}{\pi_i \pi_j} y_i y_j
$$

\end{frame}

\subsection{Design effects}

\begin{frame}\frametitle{Design-based inference, final}

Given what we just learned about means and variances in survey sampling, what are the impacts of
the survey features on analysis of survey data?

\begin{itemize}
    \item Stratification
    \item Clustering
    \item Unequal probabilities of selection
    \item Weight adjustments
\end{itemize}

Impacts could be on\ldots
\begin{itemize}
    \item Point estimates
    \item Variances --- there is so much interest in them that survey statisticians coined a special term,
        \textit{design effect}
        $$
            {\rm DEFF} = \frac{\Var[\hat\theta; \mbox{actual design}]}{\Var[\hat\theta; \mbox{SRS}]}
        $$
\end{itemize}

\end{frame}

\begin{frame}\frametitle{Stratification}

\begin{itemize}
    \item Does not have impact on $\pi_i$ and $w_i$ per se, although\ldots
    \item \ldots if different strata are sampled at different rates $\pi_i$, weights will differ
    \item Samples are taken independently between strata $\Rightarrow$ different
        expressions for the second order probabilities for units in the same vs. different strata
        \begin{itemize}
            \item $\pi_{ij} = \pi_i \pi_j; \pi_{ij} - \pi_i \pi_j=0$ 
                  so no cross-strata terms in variance
        \end{itemize}
    \item In the end of the day, stratification \textit{typically} improves precision of the estimates
        \begin{itemize}
            \item \ldots but to actually get this gain in your software output,
                you need to know what formula to use for variance
        \end{itemize}
\end{itemize}

\end{frame}

\begin{frame}\frametitle{Cluster sampling}

\begin{itemize}
    \item Does not have impact on $\pi_i$ and $w_i$ per se
    \item Units within the same cluster have \textbf{much} higher probabilities of joint selection
        than units in different clusters ($\pi_{ij} \gg \pi_i \pi_j; \pi_{ij} - \pi_i \pi_j > 0$)
    \item Variance is clearly affected, and cluster samples typically result in a loss of precision
        vis-a-vis ``flat'' samples
\end{itemize}

$$
{\rm DEFF}_{\rm cluster} = 1 + \rho_{\rm ICC}(\bar m-1)
$$
where
\begin{itemize}
    \item $\rho_{\rm ICC}$ is the \textit{intraclass correlation} $\approx$
        a fraction of total variance that is due to between-cluster variance
        \begin{itemize} \item typically low single digit \% \end{itemize}
    \item $\bar m$ is the (average) number of observations per cluster in the sample
\end{itemize}

Extreme cluster sample: systematic samples

\end{frame}

\begin{frame}\frametitle{Unequal probabilities of selection}

\begin{itemize}
    \item Probabilities of selection affect weights and hence point estimates
    \item Second order probabilities are affected $\Rightarrow$ variances are affected, as well
\end{itemize}

\medskip

(Model-based) DEFF due to unequal weighting is

$$
{\rm DEFF}_{\rm uwe} = 1 + {\rm CV}_w^2 = \frac{n \sum_{j\in\mathcal{S}} w_j^2}{\Bigl[\sum_{j\in\mathcal{S}} w_j\Bigr]^2}
$$

\medskip

\begin{itemize}
    \item The second order probabilities of selection in unequal probability sampling designs
        are one of the most convoluted problems in sampling statistics!
    \item \citet{brewer:hanif:1983} list 50 methods of unequal probability sampling, for a good measure
\end{itemize}


\end{frame}

\begin{frame}\frametitle{Weight adjustments}

\begin{itemize}
    \item Weights affect point estimates
    \item Weight adjustments are based on the sample data
        \begin{itemize} \item weights are random rather than fully-prespecified probability weights \end{itemize}
    \item Sampling variability in weights is difficult to account for
        \begin{itemize}
            \item Replicate variance estimation methods offer hope
        \end{itemize}
    \item Weight adjustments increase dispersion of weights $\Rightarrow {\rm DEFF}_{\rm uwe}$
    \item On the other hand, weight calibration improves precision of estimates
            $$
            {\rm DEFF}_{\rm calib} [\bar y] \approx 1 - R^2_{y:{\rm calib}}
            $$
            where $R^2_{y:{\rm calib}}$ is from the regression of the outcome being analyzed on the calibration variables
\end{itemize}

\medskip

\citet{deville:sarndal:1992},
\citet{henry:valliant:2015},
\citet{devaud:tille:2019}

\end{frame}

% variance estimation: linearization, replicate weights
\begin{frame}\frametitle{Variance estimation}

\begin{itemize}
    \item Analytical methods: linearization $\Rightarrow$ sandwich estimators
    \item Computational methods: replicate variance estimation
    \begin{itemize}
        \item Mimic the sampling design to create subsamples of you data
        \item Utilize variability in estimates between subsamples to inform variance estimation
        \item Balanced repeated replication (BRR; \citet{mccarthy:1969})
        \item Jackknife \citep{krewski:rao:1981}
        \item Bootstrap \citep{rao:wu:1988,sitter:1992:3bs}
        \item Successive difference (SDR)
    \end{itemize}
\end{itemize}

\bigskip

\citet{shao:1996}, \citet{rust:rao:1996}, \citet{kolenikov:2010}; 
relevant sections of \citet{heeringa:west:berglund:2017} or \citet{lumley:2010}.

\end{frame}


% methods: summary statistics, tabulations, regression models, survival models
\begin{frame}\frametitle{Summary statistics}

\begin{itemize}
    \item Starting point: estimates of totals $t_w[{\bf y}]$
        \begin{itemize} \item Estimation of proportions: define $y=0/1$ indicator \end{itemize}
    \item Point estimation: functions $g(t_w[y_1], t_w[y_2], \ldots)$
        \begin{itemize}
            \item correlation: $r(x,y) = \frac{t_w[1]t_w[xy]-t_w[x]t_w[y]}%
                {\sqrt{(t_w[1]t_w[x^2]-t_w[x]^2)(t_w[1]t_w[y^2]-t_w[y]^2)}}$
        \end{itemize}
    \item Variance estimation: sandwich
        $$
            v[g(t_w[y_1], t_w[y_2], \ldots)] = \{\nabla g\}' \, \{ v[(t_w[y_1], t_w[y_2], \ldots)] \} \, \{\nabla g \}
        $$
\end{itemize}

\end{frame}

\begin{frame}\frametitle{Regression models}

Parametric statistical models have peculiar relations with survey statistics
\citep{skinner:1989,binder:roberts:2003,binder:roberts:2009,pfeffermann:2011:whymodel}.

Census regression and normal equations:
$$
y_i = B x_i + e_i, \quad \sum_{i\in\mathcal{U}} x_i (y_i - B x_i) = 0
$$

Normal equations represent population totals (equal to zero). Plug in sample estimators:
$$
\sum_{i\in\mathcal{S}} w_j x_j (y_j - b x_j) = 0 \Rightarrow b = (X'WX)^{-1} (X'Wy), W = {\rm diag}\{ w_j \}
$$

Inference: sandwich variance estimates \citep{fuller:1975,binder:1983}.

\end{frame}

\subsection{Subpopulations / domains}

% subpopulations
\begin{frame}\frametitle{Subpopulations / domains}

\textit{Domain} $\equiv$ \textit{subpopulation} $\equiv$ a nontrivial subset of the population

\bigskip

Variance estimation complexities:
\begin{itemize}
    \item Sample size (denominator of $s^2/n$) is random even when the overall sample size is fixed by design
    \item Pairwise selection probabilities within the domain are not the true pairwise probabilities
    \item In practice, one has to work with filtering variables $d_i y_i$ instead of $y_i$ where
        $d_i$ is the domain indicator
    \item The full data set has to be used; one cannot simply drop cases that are not used for the analysis
\end{itemize}

\citet{west:berglund:heeringa:2008}

\end{frame}

\section{Other topics, one slide each}

\subsection{Software for complex survey data analysis}
% 4. Available software: R, Stata, SAS, SUDAAN. Syntax specification basics.

\begin{frame}\frametitle{Complex survey analysis software}

R | Stata | SAS | SUDAAN | SPSS

\bigskip

The bare minimum:
\begin{itemize}
    \item survey features: stratification, clustering, unequal weights
    \item methods: descriptive statistics (means, totals, tabulations)
\end{itemize}

Better working set:
\begin{itemize}
    \item survey features:
    \begin{itemize}
        \item calibrated weights
        \item replicate weights
        \item multistage samples
        \item multiphase samples
    \end{itemize}
    \item methods:
    \begin{itemize}
        \item regression models
        \item hooks for custom methods to ride on
    \end{itemize}
\end{itemize}

\end{frame}



\subsection{Advanced uses of survey data}

% 5. Advanced uses of survey data:
% small area estimation, multilevel models, multiply imputed data, survival analysis, case-control studies
\begin{frame}\frametitle{Advanced topics}

\begin{itemize}
    \item Multilevel models
    \item Small area estimation (MRP is one of the many methods)
    \item Missing (and imputed) data
    \item Case-control studies
    \item Causal inference
    \item Survival analysis
    \item Bayesian methods
\end{itemize}

\end{frame}


\subsection{Collecting survey data}
% 6. Survey data quality control: nonresponse biases, coverage biases, mode effects.

\begin{frame}\frametitle{Do you want to collect your own survey data?}

\centering

\includegraphics[scale=0.44]{tse-slide}

\end{frame}

\begin{frame}\frametitle{Focus on biases}

Statisticians are reasonably good in quantifying the sampling error. 
The ``social science'' side survey methodologists concentrate on controling 
the measurement errors part of TSE
\citep{grov:fowl:coup:lepk:sing:tour:2009,groves:lyberg:2010}.

\begin{itemize}
    \item Coverage bias
    \item Psychology of survey response
    \item Nonresponse (noncontact + refusal) bias
    \item Social desirability bias
    \item Cognitive shortcuts
    \item Mode effects
    \item Interviewer effects
    \item Multicultural, multilingual
\end{itemize}

Handbook of Health Survey Methods \citep{johnson:2015:handbook}

\end{frame}

\subsection{Resources}

\begin{frame}\frametitle{Additional resources}

Professional organizations:
\begin{itemize}
    \item Survey Research Methods Section of the ASA (SRMS)
    \item American Association for Public Opinion Research (AAPOR)
\end{itemize}

Online education:
\begin{itemize}
    \item SRMS list of degree programs: http://community.amstat.org/surveyresearchmethodssection/links
    \item Joint Program in Survey Methodology online: https://jpsmonline.umd.edu/
    \item Coursera Specialization: https://www.coursera.org/specializations/data-collection
\end{itemize}

List of competencies:
\begin{itemize}
    \item Special issue of AAPOR's \textit{Survey Practice} online journal: https://www.surveypractice.org/issue/594-vol-8-issue-2-2015
\end{itemize}

\end{frame}

\begin{frame}
    \frametitle{What I covered today}
    \tableofcontents
\end{frame}

\appendix

\section{Software}

\begin{frame}

\centering

{\Large Appendix I: Software}

\bigskip

\bigskip

Code examples: https://github.com/skolenik/ICHPS2023-svy

\end{frame}



\begin{frame}\frametitle{Complex survey analysis software}

R | Stata | SAS | SUDAAN | SPSS

\bigskip

The bare minimum:
\begin{itemize}
    \item survey features: stratification, clustering, unequal weights
    \item methods: descriptive statistics (means, totals, tabulations)
\end{itemize}

Better working set:
\begin{itemize}
    \item survey features:
    \begin{itemize}
        \item calibrated weights
        \item replicate weights
        \item multistage samples
        \item multiphase samples
    \end{itemize}
    \item methods:
    \begin{itemize}
        \item regression models
        \item hooks for custom methods to ride on
    \end{itemize}
\end{itemize}

\end{frame}


\begin{frame}\frametitle{R}

\texttt{library(survey)} --- available from CRAN \citep{lumley:2010}

\texttt{library(ReGenesees)} --- requires custom installation

\begin{enumerate}
    \item Declare your survey design (e.g., identifiers of units, strata, clusters; weight variables; etc.)
    \item Apply survey-aware functions provided by those packages to obtain design-correct inference
\end{enumerate}

Both packages support nearly any design imaginable, as well as estimation of common statistical models (e.g. GLMs)

\bigskip

Overall impression: very solid

\bigskip

\texttt{library(srvyr)} --- a tidyverse interface to \texttt{\{survey\}}

\texttt{library(sampling)} --- a suite of sampling methods \citep{tille:2006}

\end{frame}

\begin{frame}\frametitle{Stata}

Official Stata includes the suite of \texttt{svy} routines.

\begin{enumerate}
    \item \texttt{svyset} your data (identifiers of units, strata, clusters; weight variables; calibration variables)
    \item prefix nearly any estimation command with \texttt{svy} for design-correct inference
    \item clear (but complex) ways for third-party modules to comply with \texttt{svy} requirements
    \item calibrated weights (although poorly documented)
\end{enumerate}

Overall impression: very solid

Lags behind R: sampling; beats R: variety of supported models

\medskip

\citet{heeringa:west:berglund:2017}, \citet{kolenikov:pitblado:2014}.

\end{frame}

\begin{frame}\frametitle{SAS}

% TODO: include logo

The suite of \texttt{SAS PROC SVYWHATEVER} procedures:

\begin{itemize}
    \item need to declare the survey features within every procedure (copy/paste errors?)
    \item a limited set of methods coded in \texttt{SAS PROC SVYWHATEVER}
    \item other procedures produce ridiculous results with weights
\end{itemize}

\bigskip

Overall impression: OK for the general 90\% tasks

\bigskip

Lags behind R and Stata: variety of methods.

\medskip

Beats Stata with \texttt{PROC SURVEYSELECT}.

\end{frame}

\begin{frame}\frametitle{SUDAAN}

% TODO: include logo

SAS-callable custom software (SUrvey DAta ANalysis)

\begin{itemize}
    \item need to declare the survey features within every procedure (copy/paste errors?)
    \item a wider set of methods compared to SAS (including Cox regression)
    \item sophisticated weight adjustment methods
\end{itemize}

\bigskip

Overall impression: very strong within its limited scope; future development unclear

\end{frame}

\begin{frame}\frametitle{SPSS}

% TODO: include logo

\begin{itemize}
    \item weighted procedures produce ridiculous results
    \item requires a separate purchase of Complex Samples module
    \item replicate weights are not supported
    \item future development unclear
\end{itemize}

\vspace{1in}

Overall impression: meh

\end{frame}

\begin{frame}\frametitle{Software and rubrics}

\noindent
\begin{tabular}{l|ccccc}
    & R
        & Stata
            & SAS
                & SUDAAN
                    & SPSS \\
\hline
Integration
    & \texttt{library}
        & Official
            & PROC
                & SAS-
                    & Complex \\
    & \texttt{(survey)}
        & \texttt{svy:}
            & SVY*
                & callable
                    & Samples (\$\$) \\
Big three & + & + & + & + & + \\
Replicate & + & + & + & + & - \\
weights & \\
Calib $\mathbb{V}$
    & + & + & - & + & - \\
Models & ++ & ++ & + & + & ? \\
\end{tabular}

\end{frame}

\section{Advanced topics}

\begin{frame}

\centering

\Large Appendix II: Statistical models

\end{frame}



\begin{frame}\frametitle{Advanced topics}

\begin{itemize}
    \item Multilevel models
    \item Small area estimation
    \item Missing (and imputed) data
    \item Case-control studies
    \item Causal inference
    \item Survival analysis
    \item Bayesian methods
\end{itemize}

\end{frame}


\begin{frame}\frametitle{Multilevel models}

$$
y_{ij} \sim f(\theta_{ij}) \mbox{ (exponential family)}, \quad \theta_{ij} = x_{ij}'\mathbf{\beta}  + Z u_i
$$

\begin{itemize}
    \item Mixed models in biostatistics
    \item Multilevel models in social sciences
    \item If clusters and individuals are sampled in an informative way,
        variance components that are central in the maximum (quasi-)likelihood estimation
        may work in strange ways
    \item Open research question since \citet{pfef:skin:holm:gold:rasb:1998} as to how
        scale the weights at the lower levels (to tweak biases in variance estimates, efficiency gains)
\end{itemize}

\end{frame}

\begin{frame}\frametitle{Small area estimation}

Try to do more with less --- provide estimates for domains with insufficient sample size

\begin{itemize}
    \item \citet{fay:herriot:1979} model:
        $$
            g(\mbox{poverty rate}_i) = x_i'\mathbf{\beta} + \mbox{model error}_i + \mbox{sampling error}_i
        $$
    \item Unit models \citep{rao:molina:2015}
    \begin{itemize}
        \item sample model
            $$
                y_{ij} = {\bf x}_{ij}' \mathbf{\beta} + u_i + e_{ij}
            $$
        \item combined with known population totals $\Rightarrow$ estimate
            $$
                \bar Y_i = f_i \bar y_i + (1-f_i) (\bar {\mathbf X}_i ' \hat{\mathbf{\beta}} + \hat u_i)
            $$
        \item GLM-like extensions are feasible
    \end{itemize}
    \item Rediscovered independently in other literatures (economics: poverty mapping; political science: MRP),
        invariably with problems \citep{molina:rao:2010}
\end{itemize}

\end{frame}

\begin{frame}\frametitle{Imputed data}

If casewise deletion, imputation by (mean, mode, etc.), single imputation
are bad, is multiple imputation \citep{rubin:1996} a good answer?

\begin{itemize}
    \item \citet{kim:brick:fuller:kalton:2006} --- multiple imputation may not play well with the complex designs
    \item \citet{shao:sitter:1996} bootstrap:
        \begin{enumerate}
            \item Draw bootstrap sample reflecting complex design
            \item Impute once
            \item Estimate, store the pseudo-value
            \item Repeat 1--3 the ``bootstrap'' number of times (hundreds, not single digits)
            \item Combine using the bootstrap rules (although Rubin rules would be about the same)
        \end{enumerate}
\end{itemize}

\end{frame}

\begin{frame}\frametitle{Case-control studies}

\begin{itemize}
    \item Cases: sampled at 100\% rate (or close)
    \item Controls: match cases, in some appropriate sense
    \item Controls are sampled at much lower rate, $\approx$ the condition prevalence
\end{itemize}

The interest is often in the causes of the condition --- tempting to run a logistic
regression with the condition as the outcome?

\begin{itemize}
    \item Regression with weights is consistent wrt design, but likely very inefficient
    \item Regression without weights sounds scary as the design is \textbf{very} informative
    \item If sampling of the controls is not informative, then only the intercept is biased
        \citep{scott:wild:2003}
\end{itemize}

\end{frame}

\begin{frame}\frametitle{Survival analysis}

Nearly all of the expertise is encoded in Jerry Lawless' brain
\citep{lawless:2003}.

\begin{itemize}
    \item Survival analysis == instantaneous hazard rates
            $\approx$ moderately high frequency
            of observations (days); surveys $\approx$ one time operations
    \item Even in longitudinal survey studies, time between data collection
          waves ($\sim$year) may contain several events
    \item Severe left and right censoring issues
    \item Proportional hazards Cox model: \citet{binder:1992}
\end{itemize}

NCHS linked the National Death Index (NDI) with NHIS, NHANES,
some other surveys.

\end{frame}

\begin{frame}\frametitle{Bayesian methods}

Mostly forced marriage:
\begin{itemize}
    \item Bayesian methods proceed by tightening the distributions of parameters
        driven by the likelihood supplied by the data
    \item Survey inference (nonparametric!): likelihood of an SRS sample = $1/{N \choose n}$ and does not depend on $y_i$
    \item One can construct models and priors that allow reproducing some of the simple
        textbook formulae \citep{little:2012}
\end{itemize}

Bayesian methods work very well when complementing survey inference conditional on the sample
\begin{itemize}
    % \item Nonresponse analysis and adjustments
    \item Missing data imputation
    \item Small area estimation
    \item Adaptive sampling designs
\end{itemize}

\end{frame}




\section{Observational inference}

\begin{frame}\frametitle{Interface of svy and te}

Looks like we need three probability spaces:
\begin{enumerate}
	\item Superpopulation model $\xi$ to produce $(y^{(0)}_i, y^{(1)}_i, x_i)$
	\item Sampling $p$
	\item Treatment assignment $T$
\end{enumerate}


\begin{itemize}
   \item Large scale evaluation: $\xi$ creates the universe, $T$ is a randomized assignment,
           the survey follows everyone adjusting for nonresponse: $p \succ T \perp \xi$
   \item Survey experiments (e.g. question wording, mode assignment): $\xi$ creates the universe,
           $p$ draws the sample, $T$ is conditional on the sample drawn, $T \perp \xi \otimes p$. 
   \item Everything else: the probabilities spaces do not separate.
\end{itemize}

\end{frame}

\begin{frame}\frametitle{Propensity scores and weights}

\begin{itemize}
	\item \citet{dong:stuart:lenis:nguyen:2020}: 
	      literature review of the more commonly used methods to estimate PATE or PATT
			\begin{enumerate}
				\item ignore everything
				\item throw weights in as PS predictors
			   \item ignore weights in PS modeling and subclassification, 
			          use weights in outcome regression within subclasses
			   \item multiply the weights through for the outcome analysis
			   \item that is not the full factorial space! 
			\end{enumerate}
	\item \citet{ridgeway:kovalchik:griffin:kabeto:2015}: 
	      conditions on when unweighted propensity score analysis is problematic
	\item \citet{lenis:nguyen:dong:stuart:2017},
			\citet{austin:jembere:chiu:2018}: weights must be used in the outcome analysis,
			may not matter as much for PS model itself if covariates are balanced in population
			/ weighted sample 
			\begin{itemize}
					\item \citet{dong:stuart:lenis:nguyen:2020} walked it back
			\end{itemize}
	\item \citet{mercer:kreuter:keeter:stuart:2017}: 
			Rubin's framework applied to probability / non-probability samples
\end{itemize}
\end{frame}

\begin{frame}\frametitle{IPW ain't H-T}

Estimator: if the estimator is $\tau_{\rm ATE}=\frac{yt}{p}-\frac{(1-t)y}{1-p} := \sum_{i \in \mathcal{U}} \frac{y_i}{\pi_i}$, 
what are the components?

Matching the concepts from the $T$ space to the $p$ space:

\begin{itemize}
    \item Population: $\{-y^0, y^1\}$
    \item Sample: $n=1$
    \item Probabilities: $\pi_{01}=0$
    \item Variance: $(y^1+y^0)^2$
\end{itemize}

Do these make sense? 

\end{frame}


\begin{frame}[allowframebreaks]\frametitle{References}

\bibliography{everything}
\bibliographystyle{agsm}

\end{frame}



\end{document}

